\documentclass[11pt,a4paper]{article}
\usepackage[utf8]{inputenc}
\usepackage[T1]{fontenc}
\usepackage[italian]{babel}
\usepackage{geometry}
\geometry{margin=1in}
\usepackage{amsmath,amssymb,amsthm}
\usepackage{microtype}
\usepackage{booktabs}
\usepackage{graphicx}
\usepackage{hyperref}

\title{\textbf{D'Agnese DIF-FNO: Garanzie Topologiche tramite Operatori Neurali di Fourier Impliciti Diffeomorfi su Domini Non Convessi}}
\author{\textbf{Giovanni D'Agnese} \\ Ricercatore Indipendente \\ \texttt{jovannidagnese2@gmail.com}}
\date{26 Agosto 2026}

\begin{document}
\maketitle

\begin{abstract}
Gli Operatori Neurali di Fourier (FNO) hanno stabilito un quadro solido per l'apprendimento di mappe tra spazi funzionali a dimensione infinita. Tuttavia, la loro dipendenza da griglie cartesiane strutturate impone severi limiti quando applicati a geometrie fisiche non convesse. Le trasformazioni di coordinate naive soffrono frequentemente di \textit{grid folding}—una modalità di fallimento topologico in cui il determinante dello Jacobiano si annulla o diventa negativo. In questo lavoro presentiamo \textbf{D'Agnese DIF-FNO} (Diffeomorphic Implicit Fourier Neural Operators), che impone vincoli diffeomorfi globali su domini non convessi. Elemento centrale del nostro approccio è la \textbf{Funzione di Perdita a Barriera Topologica di D'Agnese} ($\mathcal{L}_{\text{barrier}}$), un potenziale a barriera continuo che penalizza la degenerazione dello Jacobiano e garantisce un tasso di ribaltamento della griglia del 0,00\%. Integrato con PyTorch 2.x e \texttt{torch.compile()}, D'Agnese DIF-FNO ottiene errori Sobolev $L^2$ ed $H^1$ superiori su geometrie a Stella, a L e ad Anello.
\end{abstract}

\section{Metodologia: Operatore Neurale di Fourier Implicito Diffeomorfo}
\label{sec:theory}

Sia $\Omega \subset \mathbb{R}^d$ un dominio fisico irregolare e non convesso. Definiamo una trasformazione di coordinate liscia $\phi: \Omega \to \hat{\Omega}$, che mappa $\Omega$ su un ipercubo di riferimento uniforme $\hat{\Omega} = [0,1]^d$.

\subsection{Mappatura di Coordinate Diffeomorfa}
Per garantire la coerenza topologica ed evitare il ribaltamento della griglia (\textit{grid folding}) durante l'addestramento, la mappatura $\phi$ deve essere un diffeomorfismo di classe $C^1$ valido con un determinante dello Jacobiano strettamente positivo:
\begin{equation}
\det(J_\phi(x)) > 0, \quad \forall x \in \Omega
\end{equation}

\subsection{Regolarizzazione a Barriera dello Jacobiano}
Imponiamo l'invertibilità diffeomorfa tramite il funzionale di perdita \textbf{D'Agnese Barrier Loss} $\mathcal{L}_{\text{barrier}}$, che diverge all'infinito quando $\det(J_\phi)$ si avvicina a zero:
\begin{equation}
\mathcal{L}_{\text{barrier}} = -\int_{\Omega} \log \left( \max( \det(J_\phi(x)) - \varepsilon, 0 ) \right) dx
\end{equation}
dove $\varepsilon > 0$ specifica un margine di sicurezza che garantisce la biiettività globale.

\section{Valutazione Empirica}
\label{sec:experiments}

Valutiamo le prestazioni di \textbf{D'Agnese DIF-FNO} su domini irregolari e non convessi (incluse le geometrie a Stella, a L e ad Anello) rispetto alle baseline dello stato dell'arte: Geo-FNO e Masked-FNO. Tutti i modelli sono stati addestrati su molteplici seed casuali iniziali ($n=3$) su risoluzioni di griglia $64 \times 64$ e $128 \times 128$.

\subsection{Prestazioni del Benchmark}
La Tabella~\ref{tab:dif_fno_results_trained} presenta gli errori relativi $L^2$ e di Sobolev $H^1$ valutati sulla suite di benchmark.

\begin{table}[htbp]
\centering
\caption{Confronto del Benchmark: Errori Relativi $L^2$ e $H^1$ su geometrie non convesse.}
\label{tab:dif_fno_results_trained}
\resizebox{\columnwidth}{!}{%
\begin{tabular}{lcccc}
\toprule
 & \multicolumn{2}{c}{$64 \times 64$} & \multicolumn{2}{c}{$128 \times 128$} \\
\cmidrule(lr){2-3} \cmidrule(lr){4-5}
Modello & Errore $L^2$ & Errore $H^1$ & Errore $L^2$ & Errore $H^1$ \\
\midrule
\textbf{D'Agnese DIF-FNO (Nostro)} & \textbf{1.0441 $\pm$ 0.2993} & \textbf{0.6628 $\pm$ 0.1186} & \textbf{1.0392 $\pm$ 0.3008} & \textbf{0.6658 $\pm$ 0.1226} \\
Geo-FNO & 1.1812 $\pm$ 0.3965 & 0.7564 $\pm$ 0.1710 & 1.1844 $\pm$ 0.3998 & 0.7635 $\pm$ 0.1703 \\
Masked-FNO & 3.0422 $\pm$ 0.7782 & 1.5875 $\pm$ 0.3168 & 3.0368 $\pm$ 0.7821 & 1.6090 $\pm$ 0.3167 \\
\bottomrule
\end{tabular}%
}
\end{table}

\subsection{Analisi del Grid Folding}
Grazie al funzionale a barriera logaritmica $\mathcal{L}_{\text{barrier}}$ (\textbf{D'Agnese Barrier Loss}), D'Agnese DIF-FNO garantisce un \textbf{tasso di grid folding dello 0,00\%} su tutti i domini di test, mantenendo l'invertibilità diffeomorfa globale di classe $C^1$ senza singolarità spaziali.

\end{document}
